\chapter{扭转 (Torsion)}

\section{受力与变形特点}

\subsection{受力特点}
外力偶作用面与杆件横截面平行。

\subsection{变形特点}
纵向线转过一角度，产生切应变（剪切角）。杆件横截面绕轴线发生相对转动。


\section{外力偶矩的计算}

\subsection{直接计算}
\begin{itemize}
    \item $M_e = F \cdot d$ （作用于盘上的力偶）
    \item $m = (F_1 - F_2) \cdot \frac{D}{2}$ （皮带传动）
\end{itemize}

\subsection{按输入功率和转速计算}
已知轴的转速为 $n$ (r/min)，输出功率为 $P$ (kW)，求外力偶矩 $M_e$：
\begin{align*}
    P &= \frac{dW}{dt} = \frac{M_e d\varphi}{dt} = M_e \omega \\
    P \times 1000 &= M_e \cdot \frac{2\pi n}{60} \\
    M_e &= \frac{60000 P}{2\pi n} \approx 9550 \frac{P}{n} \quad (\text{N}\cdot\text{m})
\end{align*}
\begin{myformula}
        M_e &= \frac{60000 P}{2\pi n} \approx 9550 \frac{P}{n} \quad (\text{N}\cdot\text{m})
    \end{myformula}
\textbf{结论：} 即低速轴更粗，高速轴更细。



\section{圆轴扭转时的应力与强度分析}

\subsection{纯剪切与薄壁圆筒}
对于薄壁圆筒（壁厚 $t \le \frac{r_0}{10}$）：
\begin{myhypothesis}{扭转假设}
    圆周线大小、间距、形状不变；纵向线倾斜同一角度。
\end{myhypothesis}
 \begin{myformula}
     \tau = \frac{T}{2\pi r_0^2 t}
\end{myformula}


\section{切应力互等定理 }

\begin{myTheorem}{切应力互等定理}
    在相互垂直的两个平面上，切应力必然成对出现，且数值相等、方向垂直于两平面的交线。其方向表现为：两个切应力要么同时指向交线，要么同时背离交线。
\end{myTheorem}


\subsection{力学本质}
该定理是由单元体的\textbf{力矩平衡条件} ($\sum M_0 = 0$) 推导得出的。
\begin{itemize}
    \item \textbf{在扭转中的体现：} 当圆轴受扭时，横截面上产生的切应力 $\tau$ 必然伴随着纵向截面（过轴线的截面）上大小相等的切应力。
    \item \textbf{破坏现象：} 对于木材等纵向抗剪能力较弱的材料，在受扭时常沿纵向裂开，这就是切应力互等定理的直观体现。
\end{itemize}

\subsection{纯剪切状态}
当单元体的各个面上只有切应力而无正应力时，称该点处于\textbf{纯剪切 (Pure Shear)} 状态。
\subsection{剪切虎克定律}
\begin{myformula}
    \tau = G\gamma, \quad G = \frac{E}{2(1+\mu)}
\end{myformula}

\subsection{圆轴扭转应力公式推导}
\begin{enumerate}
    \item \textbf{几何关系：} $\gamma_\rho = \rho \frac{d\varphi}{dx} = \rho \theta$ （切应变与半径 $\rho$ 成正比）
    \item \textbf{物理关系：} $\tau_\rho = G \gamma_\rho = G \rho \frac{d\varphi}{dx}$
    \item \textbf{静力学关系：} $T = \int_A \tau_\rho \cdot \rho \, dA = G \frac{d\varphi}{dx} \int_A \rho^2 dA = G I_p \frac{d\varphi}{dx}$
\end{enumerate}

\subsection{应力公式}
横截面上任意一点的切应力：
\begin{myformula}
    \tau_\rho = \frac{T \rho}{I_p}
\end{myformula}
横截面上的最大切应力（边缘点 $\rho = R$）：
\begin{myformula}
    \tau_{\max} = \frac{T R}{I_p} = \frac{T}{W_t}
\end{myformula}
其中 $I_p$ 为极惯性矩，$W_t = \frac{I_p}{R}$ 为抗扭截面模量。
\begin{myformula}
    I_p = \int_A \rho^2 dA, \quad W_t = \frac{I_p}{R}
\end{myformula}

\subsection{截面几何参数}
\begin{itemize}
    \item \textbf{圆轴：} $I_p = \frac{\pi d^4}{32}, \quad W_t = \frac{\pi d^3}{16} $
    \item \textbf{圆环：} $I_p = \frac{\pi (D^4 - d^4)}{32}, \quad W_t = \frac{\pi (D^4 - d^4)}{16D} $
\end{itemize}



\section{圆轴扭转变形与刚度计算}

\subsection{变形公式}
\begin{itemize}
    \item \textbf{单位长度扭转角：} 
    \begin{myformula}
        \theta = \frac{d\varphi}{dx} = \frac{T}{G I_p}
    \end{myformula}
    
    \item \textbf{相对扭转角 $\varphi$：}
    \begin{enumerate}
        \item 阶梯轴/分段等扭矩：$\varphi = \sum_{i=1}^{n} \frac{T_i l_i}{G_i I_{pi}}$
        \item 变截面或连续变扭矩：$\varphi = \int_0^l \frac{T(x)}{G I_p(x)} dx$
    \end{enumerate}
\end{itemize}
$G I_p$ 称为杆的\textbf{抗扭刚度}。

\subsection{强度条件与刚度条件}
\begin{enumerate}
    \item \textbf{强度条件：} $\tau_{\max} = \frac{T_{\max}}{W_t} \le [\tau]$
    \item \textbf{刚度条件：} $\theta_{\max} = \frac{T_{\max}}{G I_p} \cdot \frac{180}{\pi} \le [\theta] \quad (^\circ/\text{m})$
\end{enumerate}

